\documentclass[12pt, twoside]{article}
%\documentclass[12pt, twoside]{article}
\usepackage[letterpaper, margin=1in, headsep=0.2in]{geometry}
\setlength{\headheight}{0.6in}
%\usepackage[english]{babel}
\usepackage[utf8]{inputenc}
\usepackage{microtype}
\usepackage{amsmath}
\usepackage{amssymb}
%\usepackage{amsfonts}
\usepackage[nomessages]{fp} %\FPeval{\var-name}{2*sin(pi/6)}
\usepackage{siunitx} %units in math. eg 20\milli\meter
\usepackage{yhmath} % for arcs, overparenth command
\usepackage{tikz} %graphics
\usetikzlibrary{quotes, angles, arrows, arrows.meta}
\usepackage{pgfplots}
\usepgfplotslibrary{fillbetween}
\pgfplotsset{compat=1.16}
\usepackage{graphicx} %consider setting \graphicspath{{images/}}
\usepackage{parskip} %no paragraph indent
\usepackage{enumitem}
\usepackage{multicol}
\usepackage{venndiagram}

\usepackage{fancyhdr}
\pagestyle{fancy}
\fancyhf{}
\renewcommand{\headrulewidth}{0pt} % disable the underline of the header
\raggedbottom
\hfuzz=2mm %suppresses overfull box warnings

\usepackage{hyperref}

\fancyhead[LO]{La Scuola d'Italia / Dr.\ Huson / IB Math 12 \\* 29 September 2026}
\fancyhead[RO]{\fbox{\begin{minipage}{6cm}\footnotesize Nickname:\\[0.5cm]\end{minipage}}}
\fancyhead[LE]{\fbox{\begin{minipage}{6cm}\footnotesize Nickname:\\[0.5cm]\end{minipage}}}
\fancyfoot[C]{\thepage}

\begin{document}
\subsubsection*{1.5 Classwork: Box-and-Whisker Plots}

\emph{Use your GDC. Give answers exactly, or correct to three significant
figures. An outlier is a value more than $1.5 \times \text{IQR}$ below
$Q_{1}$ or more than $1.5 \times \text{IQR}$ above $Q_{3}$. On a
box-and-whisker plot, mark an outlier with a cross ($\times$).}

\begin{enumerate}[itemsep=0.15cm]

%--- Problem 1: Five-number summary from raw data; IQR; outlier test; draw the box plot [9 marks] ---
%--- Source: Original (freshly authored, AA SL 4.2), GDC One-Variable Statistics from a list ---
\item[1.]

Fifteen students record the time, in minutes, of their journey to school.

\begin{center}
$24,\; 18,\; 31,\; 12,\; 58,\; 22,\; 26,\; 15,\; 30,\; 21,\; 35,\; 20,\;
28,\; 33,\; 25$
\end{center}

\begin{enumerate}[itemsep=0.1cm]
  \item Write down the median, the lower quartile $Q_{1}$ and the upper
  quartile $Q_{3}$.
  \item Find the interquartile range.
  \item Show that $58$ is an outlier.
  \item On the scale at the bottom of the box, draw a box-and-whisker
  plot of the journey times. Mark the outlier with a cross.
\end{enumerate}

\begin{tikzpicture}
  \draw (-0.6,0) rectangle (14.6,13.4);
  \foreach \y in {5,6,...,12} \draw[dotted] (0,\y)--(14.1,\y);
  \begin{scope}[xshift=1cm, x=0.2cm]
    \foreach \t in {0,5,...,60} \draw[dotted] (\t,1.4)--(\t,4.2);
    \draw (0,1.4)--(60,1.4);
    \foreach \t in {0,5,...,60} \draw (\t,1.25)--(\t,1.55);
    \foreach \t in {0,10,...,60} \node[below] at (\t,1.25) {\footnotesize \t};
    \node at (30,0.35) {\footnotesize Journey time (minutes)};
  \end{scope}
\end{tikzpicture}

\newpage

%--- Problem 2: Read a box plot; outlier test from a five-number summary; draw and compare [10 marks] ---
%--- Source: Original (freshly authored, AA SL 4.2) ---
\item[2.]

The box-and-whisker plot at the bottom of the box shows the test scores of
the students in Class A. The scores of the students in Class B are
summarized in the table.

\begin{center}
\setlength{\tabcolsep}{10pt}\renewcommand{\arraystretch}{1.2}
\begin{tabular}{|l|c|c|c|c|c|}
\hline
Class B & Minimum & $Q_{1}$ & Median & $Q_{3}$ & Maximum \\
\hline
Score & 22 & 56 & 66 & 76 & 94 \\
\hline
\end{tabular}
\end{center}

\begin{enumerate}[itemsep=0.1cm]
  \item Write down the median and find the interquartile range of the
  Class A scores.
  \item Show that the lowest Class B score is an outlier, and that the
  highest Class B score is not.
  \item The lowest Class B score that is not an outlier is $41$. On the
  same scale as Class A, draw a box-and-whisker plot of the Class B
  scores. Mark the outlier with a cross.
  \item Compare the scores of the two classes. Make one comparison of
  the centre and one comparison of the spread, quoting values.
\end{enumerate}

\begin{tikzpicture}
  \draw (-0.6,0) rectangle (14.6,13.2);
  \foreach \y in {6,7,...,12} \draw[dotted] (0,\y)--(14.1,\y);
  \node[left] at (1.8,4.6) {\footnotesize Class A};
  \node[left] at (1.8,2.6) {\footnotesize Class B};
  \begin{scope}[xshift=-0.8cm, x=0.14cm]
    \foreach \s in {20,25,...,100} \draw[dotted] (\s,1.4)--(\s,5.4);
    \draw (20,1.4)--(100,1.4);
    \foreach \s in {20,25,...,100} \draw (\s,1.25)--(\s,1.55);
    \foreach \s in {20,30,...,100} \node[below] at (\s,1.25) {\footnotesize \s};
    \node at (60,0.35) {\footnotesize Test score};
    % Class A: 40, 52, 60, 68, 84
    \draw[thick, fill=white] (52,4.2) rectangle (68,5.0);
    \draw[thick] (60,4.2)--(60,5.0);
    \draw[thick] (40,4.6)--(52,4.6);
    \draw[thick] (68,4.6)--(84,4.6);
    \draw[thick] (40,4.35)--(40,4.85);
    \draw[thick] (84,4.35)--(84,4.85);
  \end{scope}
\end{tikzpicture}

\end{enumerate}
\end{document}
